%! suppress = Makeatletter
%! suppress = TooLargeSection
%! suppress = MissingLabel
\documentclass{article}

% Fields
\usepackage{geometry}
\geometry{top=25mm}
\geometry{bottom=35mm}
\geometry{left=20mm}
\geometry{right=20mm}
% ------------------------------------------------

% Graphics
\usepackage{color}
\usepackage{tabularx}
\usepackage{tikz}
\usepackage{blkarray}
\usepackage{graphicx}
% ------------------------------------------------

% Math
\usepackage{amsmath, amsfonts}
\usepackage{amssymb}
\usepackage{proof}
\usepackage{mathrsfs}
% Crossed-out symbols
% https://tex.stackexchange.com/questions/75525/how-to-write-crossed-out-math-in-latex
\usepackage[makeroom]{cancel}
\usepackage{mathtools}
% ------------------------------------------------

% Additional font sizes
% https://www.overleaf.com/learn/latex/Questions/How_do_I_adjust_the_font_size%3F
\usepackage{moresize}
% Additional colors
% https://www.overleaf.com/learn/latex/Using_colours_in_LaTeX
\usepackage{xcolor}
% \texttimes
\usepackage{textcomp}
% ------------------------------------------------

% Language
\usepackage[utf8] {inputenc}
\usepackage[T2A] {fontenc}
\usepackage[english, russian] {babel}
\usepackage{indentfirst, verbatim}
\usetikzlibrary{cd, babel}
% ------------------------------------------------

% Fonts
\usepackage{stmaryrd}
\usepackage{cmbright}
\usepackage{wasysym}
% ------------------------------------------------

% Code
% https://tex.stackexchange.com/questions/99475/how-to-invoke-latex-with-the-shell-escape-flag-in-texstudio-former-texmakerx
% Colored, requires --shell-escape compiling option
% \usepackage{minted}
% \setminted{xleftmargin=\parindent, autogobble, escapeinside=\#\#}
\usepackage{listings}
% ------------------------------------------------

% Custom envs
% https://tex.stackexchange.com/questions/371286/draw-a-horizontal-line-in-latex
\newenvironment{proof}{\subparagraph{\hspace{-1em}Решение:\newline}}{\par\noindent\rule{\textwidth}{0.4pt}}
% ------------------------------------------------

% Custom commands
\newcommand{\comb}[1]{\mathbf{#1}}
\newcommand{\step}{\rightsquigarrow}
\newcommand{\term}[1]{\mathbf{#1}}
\newcommand{\ap}{~}
\newcommand{\termdef}{\coloneqq}
\newcommand{\subst}[3]{\left[#2 \mapsto #3 \right] #1}
\newcommand{\eqbeta}{=_\beta}
\newcommand{\eqeta}{=_\eta}
\newcommand{\tlist}[1]{\term{[}#1\term{]}} % list-term
\renewcommand{\emph}[1]{{\color{blue} \textit{#1}}}
% ------------------------------------------------

% Head
\usepackage{fancybox,fancyhdr}
\usepackage{hyperref}
\pagestyle{fancy}
\fancyhead[R]{Студент(ка) Студентов(а)} % TODO введите ваше имя
\fancyhead[L]{ИТМО MSE, Логика 2025, Дз 10}
% ------------------------------------------------

% Numbering
% https://tex.stackexchange.com/questions/80113/hide-section-numbers-but-keep-numbering
\makeatletter
\renewcommand\thesubsection{Блок \@arabic\c@subsection.\hspace{-0.8em}}
\renewcommand\thesubsubsection{Задание \@arabic\c@subsection.\@arabic\c@subsubsection\hspace{-0.8em}}
% https://tex.stackexchange.com/questions/327689/numbering-subsubsections-with-letters
\renewcommand\theparagraph{\alph{paragraph})\hspace{-0.8em}}
% https://tex.stackexchange.com/questions/129208/numbering-paragraphs-in-latex
\setcounter{secnumdepth}{4}
\makeatother
% ------------------------------------------------

\begin{document}
    \section*{Введение}

    \subsection*{Система оценивания и общие советы}

    Оценка за курс ставится по результатам экзамена, где будет учитываться также и успеваемость по практике.
    Чтобы получить допуск на экзамен, нужно получить зачёт по практике, который ставится, если зачтены все домашние задания, кроме, может быть, одного.

    Каждое домашнее задание разделено на блоки по разновидностям задач.
    Чтобы зачесть задание нужно либо набрать локальный минимум по каждому блоку, либо набрать глобальный минимум по сумме всех блоков.
    Таким образом, если у вас какой-то блок совсем не получается и вы не можете набрать по нему локальный минимум, вы можете решить побольше в других блоках и зачесть задание, достигнув глобального минимума по общей сумме баллов решенных задач.
    Если обоими способами не получается набрать нужное количество~--- ничего страшного, потом будет достаточно возможностей добрать баллов.
    Постарайтесь тратить как можно меньше сил на тревогу и переживания.

    Ваша задача в каждом дз --- постараться как можно глубже погрузиться в как можно большее количество заданий.
    Так вы либо хорошо разберетесь сами, либо дадите себе все шансы понять разбор заданий на следующей паре так, чтобы после него уже не оставалось белых пятен.
    За неполные решения тоже ставятся баллы~--- не стесняйтесь браться и пробовать.
    Пересматривайте материалы с практик и лекций~--- там должно быть всё необходимое.
    Старайтесь как можно раньше браться за домашки, пока свежи воспоминания с практики и достаточно времени на осознание.
    Если уже хорошо знакомы с темой, не тратьте время на вводные задачки в начале блоков, начинайте с середины (но проверить базовые знания тоже может быть хорошей идеей).

    Совет: как только набираете нужное количество баллов, переключайтесь на задания по другим курсам.
    Если есть желание, решать больше можно и приветствуется --- будет бонусом.
    Но ни в коем случае не в ущерб сну и не во вред себе.

    \subsection*{Помощь}

    Не стесняйтесь обращаться за помощью к вашим преподавателям, они тут как раз для этого.
    Все контакты вы найдёте на вики.
    Нам не сложно и в радость вам помогать, а наблюдать молчаливое непонимание как раз грустно.
    Не обесценивайте свои трудности, ищите поддержки --- и найдёте.

    Также просите помощи у своих одногруппников, помогать интересно и полезно.

    Если вы помогаете, не навредите!
    Предлагайте объяснения, а не ответы.
    Отталкивайтесь от потребностей собеседника, а не от желания показать себя.
    Человеку и так может быть не очень здорово, довольно просто закопать его ещё глубже, говоря избыток умных слов и демонстрируя свою эрудицию.

    Помните: у всех разная подготовка и способности к построению определённых видов умозрительных конструкций, с этим особо ничего не сделаешь.
    Но это никак вас не определяет.
    Чем сложнее даётся, тем ценнее любой шажок.
    Не будьте слишком требовательны к себе и двигайтесь в своём темпе.
    Всё будет хорошо, гарантируем.

    Не всем просто психологически обращаться за помощью.
    Но это важно, без этого никак.
    Попробуйте относиться к этому как к части работы и постарайтесь, оно окупится.

    \subsection*{\LaTeX}

    Вам не нужно изучать \LaTeX чтобы качественно выполнить все задания.
    Достаточно просто делать по образцу того, как оформлены сами задачи.
    Однако следующие материалы могут оказаться полезными:

    \begin{itemize}
        \item \href{https://www.overleaf.com/learn}{Overleaf} --- туториалы по \LaTeX и онлайн редактор.
        \item \href{https://youtu.be/NOslSvJs29I}{LaTeX: краткое введение в качественную типографику}.
    \end{itemize}

    Репозиторий с исходниками можно скачать с \href{https://github.com/mse-fp2023/calculus-hws}{github}.

    Для работы существуют следующие варианты:
    \begin{itemize}
        \item Локальная установка \LaTeX дистрибутива и ваш любимый текстовый редактор с плагином для \LaTeX.
        \item Онлайн редактор, например, \url{https://www.overleaf.com/}.
    \end{itemize}

    Сборка:
    \begin{itemize}
        \item Если вы пользуетесь IntelliJ IDEA, конфигурация сборки добавлена в репозиторий.
        Выбираете нужную и жмакаете зелёную стрелку.
        \item Можете воспользоваться утилитой Make: вызовите \texttt{make} в корневой директории репозитория.
        \item В случае TexStudio и overleaf они собирают сами, но работать так менее удобно.
    \end{itemize}

    \begin{flushright}
        P.s.\ Добро пожаловать на курс!
        Успехов! \\
        Всецело ваша, команда курса
    \end{flushright}

    \newpage


    \section*{Дз 10. Теория вычислимости}

    Глобальный минимум: 3б.

    \subsection{Вычислимость}

    Локальный минимум: 2б.

    \subsubsection{(1б)}

    Покажите, что для всякой вычислимой функции $f$ есть такая вычислимая $g$, что область определения $g$ совпадает с областью значений $f$ и, если $f(x)$ определено, $f(g(f(x))) = f(x)$.

    \begin{proof}
        TODO % todo
    \end{proof}

    \subsubsection{(1б)}

    Приведите множество натуральных чисел, к которому $m$-сводимо любое перечислимое множество натуральных чисел.

    \begin{proof}
        TODO % todo
    \end{proof}

    \subsubsection{(1б)}

    Существует ли множество натуральных чисел, к которому m-сводимо любое другое множество натуральных чисел?

    Подсказка: подумайте про мощности множеств.

    \begin{proof}
        TODO % todo
    \end{proof}

    \subsection{Модели вычислений}

    Локальный минимум: 1б.

    \subsubsection{(1б)}

    Постройте машину Тьюринга, разворачивающую последовательность символов $\Sigma = \{0, 1\}$.

    Приложите ваше решение в виде ссылки на \url{https://turingmachinesimulator.com/}.

    \begin{proof}
        TODO % todo
    \end{proof}

\end{document}
